%! TeX Program = lualatex

\documentclass[sblidx={ancient sources, modern authors, subject}]{sblbook}

\addbibresource{sbltex.bib}

\begin{document}

\frontmatter

\title{[Title of Book]}
\author{[Your Name]}
\index{front matter!half-title page}%
\maketitle

\begin{seriespage}
  \index{front matter!series page}%
\end{seriespage}

\begin{titlepage}
  \index{front matter!title page}%
\end{titlepage}

\begin{copyrightpage}
  \index{front matter!copyright page}%
\end{copyrightpage}

\begin{dedication}
  \index{front matter!dedication}%
  If there is a dedication, it should be placed on page v.
\end{dedication}

\begin{epigraph}
  \index{front matter!epigraph}%
  If there is an Epigraph, it should be placed on page vii.
\end{epigraph}

\tableofcontents

\listoffigures

\listoftables

\begin{foreword}
\end{foreword}

\begin{preface}
\end{preface}

\begin{acknowledgements}
\end{acknowledgements}

\begin{introduction}
  \index{introduction}%
  The decision whether to number the introduction with arabic or roman
  numerals depends on the nature of the introduction and the extent to which
  it is a substantive part of the main text.

  This document models and describes how to produce a printer-ready manuscript
  for SBL Press. A printer-ready manuscript is a completely formatted and
  finished text ready to be replicated by the printer. The published work will
  look exactly like the printer-ready version of the manuscript submitted

  Much of the content of this document has been taken from
  \citetitle{prmmanual} published by SBL Press \autocite{prmmanual}.
\end{introduction}

\printbiblist[heading=biblistintoc]{abbreviations}

\mainmatter

\chapter{Page and Document Setup}

\section{Typical SBL Book Specifications: 6 × 9}

See Table \ref{bookspecs} for typical book specifications.

\begin{table}
  \centering
  \caption{SBL book specifications}
  \label{bookspecs}
  \begin{tabular}{@{}lp{5.7cm}@{}}
    \MakeUppercase{Printed margins} & .75 inch left and right, top and
    bottom\index{margins} \\
    \MakeUppercase{Usable area} & 4.5 × 7.5 inches (including
    headers\index{headers}, footers\index{footers}, and
    footnotes\index{footnotes}) \\
    \MakeUppercase{Document margins} & top: 2 inches; bottom 1.75 inches; left
    2 inches; right 2 inches (for letter paper) \\
    \MakeUppercase{Running heads} & 10/12, with 12-point spacing after, often
    with a font \emph{style} different from the main body font (e.g., italics
    or small capitals); \texttt{sblbook} uses italics\index{running heads} \\
    \MakeUppercase{Chapter titles} & 14/18, with .5-inch white space (e.g., 36
    points) between the chapter\index{chapters} title and the first paragraph
    of the main text \\
    \MakeUppercase{Body-text font/leading} & 10/12, fully
    justified\index{justification}, with .25-inch indentation of the first
    line\index{font!size} \\
    \MakeUppercase{Block quotations} & 9/11, fully justified, with .25-inch
    right and left indentation and 11 points above and below the
    paragraph\index{quotations!block} \\
    \MakeUppercase{Captions} & 9/11, center justification, with 6 points above
    and below the caption\index{captions} \\
    \MakeUppercase{Footnotes/endnotes} & 9/11, fully justified (possibly with
    a .5-point rule between the bottom of the main body text and the first
    footnote); \texttt{sblbook} uses standard \LaTeX\
    \texttt{\string\footnotesize} which is
    8/9.5\index{footnotes}\index{endnotes} \\
    \MakeUppercase{Bibliography} & 9/11, fully justified, with .25-inch
    hanging indentation, with no space between bibliographical
    entries\index{bibliography} \\
    \MakeUppercase{Indices} & 9/11, set in two columns when
    possible\index{indices} \\
    \MakeUppercase{Superscript} & 9-point font; \texttt{sblbook} uses the
    \LaTeX\ default of \texttt{\string\scriptsize} which is 7-point as 9-point
    is quite large and does not match actual SBL
    publications\index{superscripts} \\
  \end{tabular}
\end{table}

\section{Page Size}

\index{page size}%
SBL Press publishes works primarily in a 6~×~9--inch format but can
print to different dimensions, depending on the type of work and the specific
series in which it is to appear. Please consult your SBL Press editor and the
SBL Press production staff if a book you are preparing seems to require
something other than the standard 6~×~9 trim size.\index{page size!trim size}

Regardless of the trim size, allowance should be made for .75-inch
margins\index{margins} on all sides. Thus, 6~×~9 texts have a maximum usable
area\index{page size!usable area} of 4.5~×~7.5 inches; 8.5~×~11--inch texts
have a maximum usable area of 7~×~9.5 inches. The vertical dimension is the
maximum length measured from the \emph{top of the running head}\index{running
heads} to the \emph{bottom of the last line of text or footnote}.

Although most books will be printed at a 6~×~9 trim size, the files you create
should be set up on 8.5~×~11 pages.

The \texttt{sblbook} class centers the 4.5~×~7.5--inch useable
area\index{page size!usable area} on the requested paper size\index{page
size!paper size}. The default is \texttt{letterpaper}, but the class option
\texttt{a4paper} will also work.  The \texttt{trim} class option outputs a
6~×~9--inch PDF.

\section{Headers and Footers}

\index{headers}\index{footers}\index{running heads}%
The \emph{header} (or running head) is the text that appears on each page
above the main text. Except for pages that begin a new chapter or section, all
page numbers should appear in the header at the outside margin. The book or
part title should appear centered in the header of verso\index{verso}
(left-hand/even) pages, and the chapter or part title should appear in
recto\index{recto} (right-hand/odd) page headers, except on the pages where
the titles first appear. In multiauthored works, contributors’ names should
appear on verso\index{verso} pages and chapter names on recto\index{recto}
pages. It is advisable to use a font style different from the body font in
page headers (e.g., italics or small capitals). The format should follow the
size and leading of the main text (typically 10/12), with 12-point spacing
after the text.

SBL Press books typically include a \emph{footer}\index{footer} (or
folio) at the bottom of a page only for the page number on the first page of a
chapter or section. A footer should include only the page number (not a book
or part title) and should be centered on the page.  The format should follow
the size\index{font!size} and leading\index{font!leading} of the main text
(typically 10/12), with 12-point spacing before the text.

\section{Chapter Title Pages}

\index{chapters!title pages}%
The first page of a chapter should begin approximately one-third down from the
top of the page. The chapter heading should include the chapter number (if
applicable, as a numeral, not spelled out) and the chapter title itself on the
following line. It is preferable to leave approximately one-half inch of
“white space” between the chapter title and the first paragraph of the main
text. As noted above, chapter title pages do not include a running head but
rather a footer with the page number centered below the text.

\section{Pagination}

The first four pages of the book (i–iv) are referred to as the \emph{front
matter}\index{front matter} (the half-title page\index{front matter!half-title
page}, the series page\index{front matter!series page}, the title
page\index{front matter!title page}, and the copyright page\index{front
matter!copyright page}) and will be prepared by SBL Press production staff in
accordance with the series format. Page v is the first page for which the
author is responsible. If there is a dedication\index{front
matter!dedication}, it should be placed on page v and the table of contents on
page vii, page vi being blank. Remember: the table of contents lists only the
items that follow it. Table \ref{frontmattersequence} offers the typical
sequence of a book’s front matter.

\begin{table}
  \centering
  \caption{Sequence of a book's front matter}
  \label{frontmattersequence}
  \begin{tabular}{lr}
    Half-title page & i \\
    Series page & ii \\
    Title page & iii \\
    Copyright page & iv \\
    Dedication (if applicable) & v \\
    blank (if applicable) & vi \\
    Table of contents & vii \\
    Foreword (optional) & recto\index{recto} \\
    Preface (optional) & recto\index{recto} \\
    Acknowledgments (optional) & recto\index{recto} \\
    Introduction (if not part of text) & recto\index{recto} \\
    Abbreviations & recto\index{recto} \\
  \end{tabular}
\end{table}

\chapter{Font Selection}

\section{Style and Size}

Texts should be prepared in a good quality, proportionally spaced, serif font,
such as (New) Baskerville\index{font!Baskerville}\index{font!New Baskerville},
Times\index{font!Times} (New) Roman\index{font!Times New Roman}, New Century
Schoolbook\index{font!New Century Schoolbook}, or
Palatino\index{font!Palatino}. SBL Press books typically use a 10-point font
for the body text and an 8- or 9-point font for footnotes or endnotes.

The \texttt{sblbook} class uses the Baskervaldx\index{font!Baskervaldx} font
by default, but an alternative can be specified by passing options to the
\textsf{sblfonts} package using the \texttt{sblfonts} class option. See the
documentation for the \textsf{sblfonts} package for details.

Chapter titles are usually set in a 14-point font. It is preferable to use the
same font face throughout the book for the main body text, running
heads,\index{running heads} headings, subheadings, block quotations,
footnotes, bibliography, and indices.\index{indices}

If available, \texttt{sblbook} will use the SBL fonts for Greek and Hebrew
with Unicode engines otherwise the class falls back to Alegreya for Greek and
Frank Ruehl CLM for Hebrew. The class sets up the \texttt{babel} package for
\texttt{polytonicgreek} and \texttt{hebrew}. Other languages can be loaded
manually or on the fly with \texttt{lualatex}.

\section{Headings and Subheadings}

The design of headings and subheadings is left to the discretion of the
author. As a rule (and as in this manual), first-level headings are centered,
and second-level headings are justified left, but other consistent styles are
acceptable. Generally it is best also to distinguish headings and subheadings
from the body text through use of a different font style (e.g., small capitals
with title-case capitalization; small capitals with no initial capitalization;
italics). Finally, headings that are followed by \emph{less} than two lines of
text before the end of a page should be avoided.\index{widows}

\section*{This Is an Example of an A-Level Heading}

\subsection*{This Is an Example of a B-Level Subheading}

\subsubsection*{This Is an Example of a C-Level Subheading.}

When you use only two levels of heading, they can be centered and flush left,
respectively, both with white space before and after the heading (settings
here: 12-point spacing before and after A-level heading; 6-point spacing
before and after B-level subheading). If you use a third level of heading, the
paragraph should have space before it, and the following text should be run
in, as in this paragraph.

\section{Paragraphs}

First lines of paragraphs are normally indented one-quarter inch. The first
line of a paragraph following a chapter title, heading, or subheading (e.g.,
this paragraph) should not be indented.

Generally the text in books is \emph{fully justified}---each line is exactly
the same length. Occasionally a line will have an unsightly space (e.g., when
the line ends in a word that cannot be divided).  This problem can be
corrected by a minor rewrite.

\subsection{Block Quotations/Extracts}

A quotation of five or more lines should be set off as a separate paragraph from
the main text and formatted differently from the body text.
\begin{quote}
  This paragraph is an example of a block quote. The margins for block
  quotations should be set in one-quarter inch from the right and left margins
  of the body text, and it is customary to use a font size 1 point smaller
  than the body font.
\end{quote}
Finally, it is preferable to leave some white space above and below the
extract (6 to 11 points). Thus, typical specifications for a block quote would
be a 9-point font with 11-point leading (9/11), indent of .25 inch right and
left, and 11 points above and below the paragraph.

\section{Footnotes and Endnotes}

\index{footnotes}%
To provide the greatest ease of use for readers, SBL Press prefers footnotes
to endnotes\index{endnotes}. Footnote numbers should restart their count at
one at the beginning of each chapter. Footnote text should be set in a font
size smaller than the main body text, usually 8 or 9 point, and leading should
be adjusted accordingly.\footnote{This is an example
footnote\index{footnotes}. Note that SBL Press prefers the first line of
footnotes to be indented. See
\url{https://sblhs2.com/2016/06/21/indenting-footnotes/}.} Thus, a typical
font/leading specification for footnotes is 8/10 or 9/11. It is acceptable but
not mandatory to set a thin rule (line) between the end of the main text and
the first footnote. If you insert such a rule, it is typically .5 point in
width, begins at the left margin, and extends approximately 1.5 inches to the
right.

The note number should be superscripted in the main text. In the footnote, the
note number should match the rest of the footnote. The typical specifications
for a superscripted number is 9-point font raised 3
points.\footnote{The specification of a 9-point font raised 3 points for
superscripts is quite large and not what the SBL Press does in any of their
actual published books. I've kept the \LaTeX\ default which is
\texttt{\string\scriptsize} (7-point) raised 3.5 point.}

\subsection{Bibliographies and Indices}

To conserve space, bibliographies and indices should be set in a font size
smaller than the main body text. In most cases, then, bibliographies and
indices will be set in 9-point font with 11-point leading (9/11). Indices
should be set in at least two columns.\index{indices}

The \texttt{sblbook} class supports a \emph{Subject Index}, \emph{Modern
Authors Index} and \emph{Ancient Sources Index}. Each index can be enabled
independently by passing one or more of \texttt{ancient sources},
\texttt{modern authors} and \texttt{subject} to the \texttt{sblidx} class
option. Modern authors from \texttt{biblatex-sbl} citations are automatically
added to the \emph{Modern Authors Index} when cited. Scriptural references
cited using \verb|\pibibleverse| from the \texttt{bibleref-parse} package are
automatically added to the \emph{Ancient Sources Index}. Each of the three
indices is formatted in a slightly different way inline with the requirements
of the instructions on preparing indices at
\url{https://www.sbl-site.org/assets/pdfs/pubs/Indexing_SBL.pdf}.

\section{Orphans and Widows}

If at all possible, avoid \emph{orphans}\index{orphans} and
\emph{widows}\index{widows} when typesetting your manuscript. Orphans and
widows are single lines that become separated from the rest of their paragraph
because of a page break. New paragraphs should not begin at the bottom of a
page if there is room for only one line, and the final line of a paragraph
should not stand alone at the top of a page.

\section{Stacking}

It is also wise to minimize “stacking”\index{stacking} in the text. Stacking
occurs when two consecutive lines begin or end with the same word(s) or two
consecutive lines end with a hyphen. Ths class will not place two hyphens on
consecutive lines, but it's best to rewrite paragraphs where consecutive lines
begin or end with the same word. You can highlight them (along with rivers) by
placing
\texttt{\string\PassOptionsToPackage\{draft,\,homeoarchy,\,rivers\}\{impnattypo\}}
just before \verb|\documentclass{sblbook}| (the \texttt{luatex} engine is
required).

\chapter{Citations}

\index{citations}%
Citations should be referenced using \verb+\autocite+ or \verb+\autocites+ for
single volume resources and \verb+\avolcite+ or \verb+\avolcites+ for
multi-volume resources. These macros will place the citation in a footnote.
\verb+\parencite+, \verb+\parencites+, \verb+\pvolcite+, and \verb+\pvolcites+
place citations in parentheses.

For example, citing a classical primary source \ptranscite{tacitus:ann};
citing a lexicon \autocite[\foreignlanguage{polytonicgreek}{παρρησία}]{BDAG};
citing volume 2 of a multi-volume commentary
\avolcite{2}[125]{dahood:1965-1970}; citing two resources by the same author
in a single footnote \autocites[504]{harrington:1970}[241]{harrington:1986};
citing a resource for the second time \autocite[505]{harrington:1970}; citing
an article in an edited collection \autocite{collins:1986}; citing another
article in the same edited collection \autocite{attridge:1986}. Citing an
author in the main body of the text can be done using \verb+\citeauthor+,
e.g., \citeauthor{scott+etal:1993} and \citeauthor{robinson+koester:1971}
\autocite[See][]{scott+etal:1993,robinson+koester:1971}.

\chapter{Greek and Hebrew}

I strongly recommend using the \texttt{luatex} engine for multilingual work,
especially for right-to-left languages. Inline Greek and Hebrew need to be
marked up with \texttt{xetex} and \texttt{pdftex}. With \texttt{pdftex} Hebrew
pointing support is limited and cantillation marks are not supported at all.
\begin{quotation}
  \selectlanguage{polytonicgreek}
  Ἐν ἀρχῇ ἦν ὁ λόγος, καὶ ὁ λόγος ἦν πρὸς τὸν θεόν, καὶ θεὸς ἦν ὁ λόγος.
  οὗτος ἦν ἐν ἀρχῇ πρὸς τὸν θεόν.
  \foreignlanguage{american}{(\pibibleverse{John 1:1})}
\end{quotation}
\begin{quotation}
  \selectlanguage{hebrew}
  בּראשׁית בּרא אלהים את השּׁמים ואת הארץ׃ והארץ היתה תהוּ ובהוּ וחשׁך על־פּני תהום
  ורוּח אלהים מרחפת על־פּני המּים׃
  \foreignlanguage{american}{(\pibibleverse{Genesis 1:1-2})}
\end{quotation}
Inline Greek (\foreignlanguage{polytonicgreek}{Ἐν ἀρχῇ}) and Hebrew
(\foreignlanguage{hebrew}{בּראשׁית}) are also possible.\footnote{Bible
references in footnotes, e.g., \pibibleverse{Deuteronomy 6:4} and
\pibibleverse{1 Peter 3:18} should appear in the index with their note
number.}

\appendix

\chapter{Appendix Title}

\index{appendices}%
Each appendix should have a number and a title, unless there is only one
appendix, in which case the appendix does not need a number. Every appendix
requires a heading, so if you are including a preexisting document you will
need to type a heading (i.e., the appendix number and title) on that document
so that it conforms to your numbered appendices.

An appendix is formatted like the first page of a chapter. Locate page numbers
at the bottom center of the first page of each appendix and at the top right
corner of subsequent pages. If the appendix is already numbered, put those
page numbers in square brackets. Page numbering for the appendices is
consecutive with the rest of the book.

Margins for the appendices should be the same as the rest of the book. You
may need to reduce the content of the appendix to fit the margins.

\backmatter

\nocite{NIDNTT, Jastrow, DMBI, mclay:2006, oday:intertextuality, rad:1990}

\printbibliography[heading=bibintoc]

\index{block quotations|see{quotations}}
\index{copyright page|see{front matter}}
\index{dedication|see{front matter}}
\index{epigraph|see{front matter}}
\index{folio|see{footer}}
\index{half-title page|see{front matter}}
\index{series page|see{front matter}}
\index{title page|see{front matter}}

\index{headers|seealso{running heads}}
\index{running heads|seealso{headers}}
\index{recto|seealso{verso}}
\index{verso|seealso{recto}}

\SBLPrintIndices

\end{document}

